\chapter{绪论}


\section{研究背景与意义}

\zhdigits*{\the\year} 年\zhnumber{\the\month} 月\zhnumber{\the\day} 日\;\xxivtime
随着汽车产业、智能移动设备以及基于位置服务的快速发展，带有地理位置信息的时空数据呈现爆炸式增长。其中，路网数据和轨迹数据是两个代表性的时空数据。路径规划、轨迹分析等基于时空数据挖掘的位置服务在城市计算中扮演着日益重要的角色。一个动态路网是由顶点集和边集构成的，其中顶点代表路段的连接点，交叉点；连边则代表具体的路段。每条边连接着两个顶点。为了便于描述，我们假定所有行程请求的起点和终点位置都已经被映射到离它们最近的顶点上。

路径规划不仅直接影响人们的日常出行方式，也关系到物流运输、智慧城市管理以及能源优化等多个领域。近年来，随着网络购物、共享出行和自动驾驶技术的发展，城市路网中行程请求呈现出高频、密集和动态的特点，传统基于单次行程或静态路径集合的最优路径规划方法已难以适应实时、连续的行程请求流。

综上，城市计算中面临的核心挑战包括大规模数据处理、动态环境适应以及对象间潜在交互的准确建模。围绕路径规划、联合移动模式挖掘、接触搜索与轨迹相似连接等任务，设计高效算法与索引结构不仅具有理论价值，也具有显著应用潜力：一方面可提升城市交通系统的调度效率与拥堵缓解能力；另一方面可支持公共卫生管理、生态监测及智能出行等多领域的决策优化。基于此，本文旨在系统研究城市计算中高效算法设计问题，通过统一的时空数据建模与高性能计算框架，实现对移动对象行为的实时、准确与可扩展分析，为智慧城市应用提供技术支撑。



\section{研究目标}

\section{研究所面临的挑战}

\section{研究内容}
在此背景下，本研究聚焦于四类关键问题：持续最优路径规划、联合移动模式挖掘、接触搜索与轨迹相似连接。这些问题均涉及海量移动对象和复杂时空约束，典型研究意义如下：

\subsection{面向大规模行程请求流的持续最优路线组合规划} 
交通拥堵是城市化进程中普遍存在的问题。传统路径规划往往基于单一行程或当前交通状态进行局部优化，忽略了新到达请求对未来路网负载的影响。面对持续到达的行程请求流，单次最优路径规划可能导致集中拥堵，降低整体通行效率。因此，需要设计面向持续请求流的全局协同路径规划算法，实现对时间窗口内行程请求的联合优化，从整体视角缓解交通拥堵，提高路网利用率。

\subsection{基于轨迹数据的群体联合移动模式挖掘} 
在大规模移动对象环境中，识别群体运动模式（如 flock、convoy、swarm 等）可揭示潜在的协同行为与社会动态规律。这类模式在拼车推荐、交通拥堵预测、异常行为检测以及生态监测中具有广泛应用。高效挖掘联合移动模式的核心挑战在于对象轨迹在时间上可能暂时离群或不连续，导致候选群体数量指数级增长，需要设计高效的索引结构与过滤策略以保证算法可扩展性。

\subsection{大规模移动对象上的持续密切接触发现} 
在公共卫生管理与疫情防控中，及时识别个体间潜在接触事件对于控制疾病传播具有关键意义。由于移动对象的轨迹数据通常存在采样间隔与定位误差，仅依赖离散采样点判断接触关系容易产生误判。面向大规模连续轨迹数据的接触搜索算法必须支持高并发、低延迟处理，并在动态更新环境下保证结果的准确性与一致性，为公共卫生决策提供数据支撑。

\subsection{面向时空运动范围的轨迹相似性连接方法} 
传统轨迹相似性连接方法多依赖离散采样点计算对象间距离，难以刻画未观测时间区间内的潜在交互。IS-Join 通过引入运动范围的概念，将对象在给定时间段内可能经过的空间区域作为分析单元，采用交集相似性度量量化对象间潜在重叠关系。这一方法可应用于交通监控、野生动物迁徙分析、接触追踪等场景，在稀疏采样或不确定条件下仍能稳健地识别潜在关联。


\section{论文的主要贡献}

\section{论文的组织结构}
